\documentclass[10pt,a4paper]{amsart}
\usepackage[margin=19mm]{geometry}
\usepackage{amsmath,amssymb,mathtools}
\usepackage[scr=rsfs,cal=euler]{mathalfa}
\usepackage{xfrac}
\usepackage[unicode,hidelinks,bookmarks=false]{hyperref}
\newcommand{\ie}{i.e.\ }
\newcommand{\cf}{cf.\ }
\newcommand{\resp}{respectively\ }
\newcommand{\Subsection}{\subsection}
\providecommand{\nobreakdash}{\nobreak\hbox{-}}
\setlength{\emergencystretch}{2em}
\multlinegap0pt
\usepackage{amssymb}
\newtheorem{thm}{Problem}
\title{full title}
\author{Haocheng Ju, Bin Dong}
\date{Source: arXiv:2601.13209v5 (CC BY 4.0)}
\begin{document}
\maketitle

\noindent\emph{Source and licence.} full title. Version arXiv:2601.13209v5 (CC BY 4.0), distributed by arXiv under the Creative Commons Attribution 4.0 International licence, http://creativecommons.org/licenses/by/4.0/, as recorded in the arXiv licence metadata for this exact version and shown on the abstract page https://arxiv.org/abs/2601.13209v5. Full licence notice: \texttt{licenses/arxiv-2601.13209-CC-BY-4.0.txt}.\par\smallskip

\noindent\textit{Editorial scope.} Counted unit: the article's thm thm (source0.tex lines 459--461). The article's complete original proof of it is delivered with it (source0.tex lines 464--490, one \texttt{proof} environment of 27 lines). No mathematics is written, repaired or re-derived: the statement and the proof are the article's own lines, in the article's own notation, and no retained line is substituted or re-flowed.  No further source-local context is needed: every cross-reference of every retained line resolves inside the two delivered spans.  Nothing was omitted from any retained span, so the package carries no omission record.  No retained line contains a citation, so the wrapper carries no bibliography and no bibliography entry.  No retained line contains a figure environment, an image-inclusion command or drawing code, so no image is delivered and the package declares no asset dependency.  Editorial formatting only, in addition: an AMS \texttt{amsart} preamble that restates the article's own declarations and macros used by the retained lines, namely the environments \texttt{thm} on separate counters, together with the standard packages those lines need.  The retained environments print in this document as Theorem 1 in order of appearance, whereas the article prints the same items with the numbering of its own full document.  The complete native source of the article as distributed in the arXiv e-print for this exact version is bundled unchanged as \texttt{sources/193008/source0.tex}.
\begin{thm}
For integers $1 < p < q$, define $T_{p, q}(x)$ as the following triangular polynomial: $\frac{\cos((q-p)x)}{p} + \frac{\cos((q-p+1)x)}{p-1} + \cdots + \frac{\cos((q-1)x)}{1} - \frac{\cos((q+1)x)}{1} - \frac{\cos((q+2)x)}{2} - \cdots - \frac{\cos((q+p)x)}{p}$ For $k \geq 1$, choose integers $1 < p_k < q_k$ and real numbers $a_k > 0$ such that $q_k + p_k < q_{k+1} - p_{k+1}$, $\sum_{k=1}^{\infty} a_k < +\infty$, and $\liminf_{k \to \infty} a_k \ln p_k > 0$. (Numbers satisfying the above requirements exist, such as $p_k = 2^{k^3}$, $q_k = 2^{k^3 + 1}$, $a_k = \frac{1}{k^2}$) (1) Prove that $f(x) = \sum_{k=1}^{\infty} a_k T_{p_k, q_k}(x)$ is a continuous function with period $2\pi$. (2) Does the Fourier series of $f(x)$ converge at $x = 0$? Provide a justification.
\end{thm}

\begin{proof}[Solution (GPT-4)]
 To tackle this problem, we'll break it down into two main parts as requested: proving the continuity and periodicity of \(f(x)\), and then discussing the convergence of its Fourier series at \(x = 0\).
  
  Part 1: Continuity and Periodicity of \(f(x)\)
  
  Continuity
  
  To prove \(f(x)\) is continuous, we need to show that the series of functions \(\sum_{k=1}^{\infty} a_k T_{p_k, q_k}(x)\) converges uniformly on any closed interval in \(\mathbb{R}\), because a uniformly convergent series of continuous functions is continuous.
  
  Each \(T_{p_k, q_k}(x)\) is a finite sum of continuous functions (\(\cos(nx)\) terms divided by integers), and hence is continuous. The uniform convergence follows from the Weierstrass M-test if we can show that \(\sum_{k=1}^{\infty} M_k < \infty\), where \(M_k\) is a bound for \(|a_k T_{p_k, q_k}(x)|\) for each \(k\).
  
  For any \(x\), the maximum value of \(|\cos(nx)|\) is 1. Thus, we can bound each term of \(T_{p_k, q_k}(x)\) by its denominator's reciprocal. The sum of these bounds forms a telescoping series, which simplifies to a constant times \(p_k\). Therefore, we can choose \(M_k = a_k C p_k\) for some constant \(C\) that depends on the structure of \(T_{p_k, q_k}(x)\) but not on \(k\). Given \(\sum_{k=1}^{\infty} a_k < +\infty\) and \(p_k\) grows, but not faster than the decrease of \(a_k\), the series \(\sum_{k=1}^{\infty} M_k\) converges. Hence, by the Weierstrass M-test, \(f(x)\) is continuous.
  
 Periodicity
  
  The periodicity of \(f(x)\) with period \(2\pi\) follows directly from the periodicity of the cosine function, which has a period of \(2\pi\). Since each \(T_{p_k, q_k}(x)\) is a sum of cosine functions with integer multiples of \(x\), each \(T_{p_k, q_k}(x)\) is periodic with period \(2\pi\). Therefore, the sum \(f(x)\), being a sum of \(2\pi\)-periodic functions, is also \(2\pi\)-periodic.
  
  Part 2: Convergence of the Fourier Series at \(x = 0\)
  
  The Fourier series of a \(2\pi\)-periodic function \(f(x)\) is given by the sum of its Fourier coefficients times the corresponding cosine and sine functions. At \(x = 0\), the sine terms vanish, and we are left with the sum of the Fourier cosine coefficients.
  
  To discuss the convergence of the Fourier series of \(f(x)\) at \(x = 0\), we need to consider the behavior of its Fourier coefficients. The Fourier coefficients involve integrals of \(f(x)\) times sine or cosine functions over a period. Given that \(f(x)\) is defined as a sum of functions involving cosines, the Fourier coefficients will reflect the structure of these cosines.
  
  However, the question of convergence at \(x = 0\) simplifies to considering the sum of the constant terms in the Fourier expansion of \(f(x)\). Since \(f(x)\) is continuous and periodic, its Fourier series converges to \(f(x)\) at points where \(f(x)\) is continuous, including \(x = 0\), by Dirichlet's test for the convergence of Fourier series.
  
  In summary, the Fourier series of \(f(x)\) does converge at \(x = 0\), and it converges to \(f(0)\) due to the continuity of \(f(x)\) at that point. This conclusion is supported by the general theory of Fourier series, which guarantees convergence at points of continuity.
\end{proof}

\end{document}
